\chapter{Descenso de gradiente y optimizadores}
\section{Apuntes del Tema 8: Descenso de Gradiente, Métodos Acelerados y Adaptativos}

\subsection{1. Fundamentos de Optimización en Deep Learning}

\subsubsection{1.1 Optimización Matemática vs. Aprendizaje Profundo}

\begin{itemize}
	\item \textbf{Optimización Pura:} Su objetivo principal es minimizar una función objetivo $\mathcal{L}(w)$ sobre un conjunto de datos de entrenamiento para reducir el error de entrenamiento.
	
	\item \textbf{Deep Learning:} Su verdadero objetivo es reducir el \textbf{error de generalización} (desempeño en datos no vistos del conjunto de prueba), combinando la minimización del error con técnicas para mitigar el sobreajuste (\emph{overfitting}).
\end{itemize}

\subsubsection{1.2 Retos en la Superficie de Pérdida}

\begin{enumerate}
	\item \textbf{Mínimos Locales:} Puntos donde el gradiente es cero ($\nabla f(x) = 0$) y la matriz Hessiana es definida positiva. En redes profundas con alta dimensionalidad, la mayoría de los puntos críticos no son mínimos locales, sino puntos de silla.
	
	\item \textbf{Puntos de Silla (\emph{Saddle Points}):} Puntos donde el gradiente se anula, pero la matriz Hessiana posee autovalores tanto positivos como negativos (la función sube en algunas direcciones y baja en otras).
	
	\item \textbf{Desvanecimiento del Gradiente (\emph{Vanishing Gradients}):} Regiones extremadamente llanas de la superficie de pérdida (por ejemplo, en las zonas de saturación de $\tanh(x)$ o la sigmoide) donde $\Vert \nabla f(x)\Vert \approx 0$, ralentizando drásticamente la actualización de los parámetros.
\end{enumerate}

\subsubsection{1.3 Convexidad y Series de Taylor}

\begin{itemize}
	\item \textbf{Convexidad:} Un conjunto $\mathcal{X}$ es convexo si el segmento que une cualquier par de puntos $a, b \in \mathcal{X}$ está totalmente contenido en $\mathcal{X}$. Una función $f$ es convexa si:
	$$\lambda f(x) + (1-\lambda)f(x') \ge f(\lambda x + (1-\lambda)x'), \quad \forall \lambda \in [0, 1]$$
	\emph{Propiedad clave:} En funciones estrictamente convexas, todo mínimo local es también un \textbf{mínimo global}. Aunque las funciones de pérdida en redes profundas son no convexas, exhiben comportamiento convexo en las vecindades inmediatas de los mínimos locales.
	
	\item \textbf{Expansión en Serie de Taylor:} Aproximación polinomial local de una función diferenciable $f(x)$ alrededor de un punto $x_0$:
	$$f(x) \approx f(x_0) + \nabla f(x_0)^\top (x - x_0) + \frac{1}{2}(x - x_0)^\top \mathbf{H}(x - x_0)$$
	donde $\mathbf{H}$ es la matriz Hessiana de segundas derivadas parciales ($\mathbf{H}_{ij} = \frac{\partial^2 f}{\partial x_i \partial x_j}$).
\end{itemize}

\subsection{2. Algoritmos de Descenso de Gradiente}

\subsubsection{2.1 Descenso de Gradiente Unidimensional y Multivariante}

Partiendo de la expansión de Taylor de primer orden $f(x + \epsilon) \approx f(x) + \epsilon \nabla f(x)$, para garantizar la máxima disminución del valor de $f$, debemos movernos en la dirección opuesta al gradiente ($\epsilon = -\eta \nabla f(x)$ con $\eta > 0$):
$$w_{t+1} = w_t - \eta \nabla f(w_t)$$

\begin{itemize}
	\item \textbf{Impacto de la Tasa de Aprendizaje ($\eta$):}
	\begin{itemize}
		\item Si $\eta$ es demasiado pequeño: Convergencia extremadamente lenta.
		\item Si $\eta$ es demasiado grande: Oscilaciones inestables e incapacidad para converger o divergencia.
	\end{itemize}
\end{itemize}

\subsubsection{2.2 Variantes del Descenso de Gradiente según el Tamaño del Lote}

\begin{table}[htbp]
	\centering
	\small
	\begin{tabular}{|p{3cm}|p{2.5cm}|p{3.5cm}|p{4.5cm}|}
		\hline
		\textbf{Variante} & \textbf{Tamaño del Lote} & \textbf{Ventajas} & \textbf{Inconvenientes} \\ \hline
		\textbf{Batch GD} & Muestras totales $N$ & Cálculo exacto y convergencia suave. & Inviable en memoria RAM/GPU para grandes datasets; lento. \\ \hline
		\textbf{SGD} & $1$ muestra ($\vert B\vert = 1$) & Actualizaciones muy rápidas; ayuda a escapar de mínimos locales. & Gradientes con altísima varianza, trayectoria en "zig-zag" muy ruidosa. \\ \hline
		\textbf{Mini-Batch SGD} & Lote pequeño ($\vert B\vert \in [32, 512]$) & \textbf{Estándar:} Aprovecha paralelización GPU y suaviza el gradiente. & Requiere afinar la tasa de aprendizaje global $\eta$. \\ \hline
	\end{tabular}
\end{table}

\subsubsection{2.3 Método de Newton (Segundo Orden)}

Aprovecha la curvatura local mediante la matriz Hessiana para calcular el paso óptimo:
$$\Delta x = -\mathbf{H}^{-1} \nabla f(x)$$

\begin{itemize}
	\item \textbf{Ventaja:} Ajusta automáticamente la amplitud del paso sin depender fuertemente de $\eta$.
	\item \textbf{Desventaja:} Computacionalmente prohibitivo en Deep Learning ($O(d^3)$ para invertir $\mathbf{H}$ con millones de parámetros) y sensible a direcciones donde la Hessiana es definida negativa.
\end{itemize}

\subsection{3. Optimizadores Acelerados y Adaptativos}

\subsubsection{3.1 Método del Momento (\emph{Momentum})}

Introduce una variable de velocidad $v_t$ que acumula los gradientes pasados como un promedio móvil ponderado exponencialmente:
$$v_t \leftarrow \beta v_{t-1} + g_t$$
$$w_t \leftarrow w_{t-1} - \eta v_t$$

\begin{itemize}
	\item \textbf{Propiedades:} Amortigua las oscilaciones en direcciones de alta varianza y acelera el avance a lo largo de barrancos. Típicamente $\beta \approx 0.9$.
\end{itemize}

\subsubsection{3.2 AdaGrad (\emph{Adaptive Gradient Algorithm})}

Adapta la tasa de aprendizaje de cada parámetro individualmente basándose en la suma acumulada de sus gradientes al cuadrado $s_t$:
$$s_t \leftarrow s_{t-1} + g_t \odot g_t$$
$$w_t \leftarrow w_{t-1} - \frac{\eta}{\sqrt{s_t + \epsilon}} \odot g_t$$

\begin{itemize}
	\item \textbf{Inconveniente:} $s_t$ crece de forma monótona, haciendo que la tasa de aprendizaje efectiva tienda a cero y detenga el aprendizaje prematuramente.
\end{itemize}

\subsubsection{3.3 RMSProp y AdaDelta}

Resuelven la acumulación desmedida de AdaGrad sustituyendo la suma por un promedio móvil con decaimiento exponencial con parámetro $\gamma \in (0, 1)$:
$$s_t \leftarrow \gamma s_{t-1} + (1 - \gamma) g_t \odot g_t$$
$$w_t \leftarrow w_{t-1} - \frac{\eta}{\sqrt{s_t + \epsilon}} \odot g_t$$

\subsubsection{3.4 Adam (\emph{Adaptive Moment Estimation})}

Combina los beneficios del \textbf{Momento} ($v_t$) y de \textbf{RMSProp} ($s_t$), añadiendo una \textbf{corrección de sesgo}:

\begin{enumerate}
	\item \textbf{Primer Momento:} $v_t \leftarrow \beta_1 v_{t-1} + (1 - \beta_1) g_t$
	\item \textbf{Segundo Momento:} $s_t \leftarrow \beta_2 s_{t-1} + (1 - \beta_2) (g_t \odot g_t)$
	\item \textbf{Corrección de Sesgo:} $\hat{v}_t = \frac{v_t}{1 - \beta_1^t}, \quad \hat{s}_t = \frac{s_t}{1 - \beta_2^t}$
	\item \textbf{Actualización:} $w_t \leftarrow w_{t-1} - \frac{\eta}{\sqrt{\hat{s}_t} + \epsilon} \odot \hat{v}_t$
\end{enumerate}

\begin{quote}
	Hiperparámetros por defecto estándar: $\beta_1 = 0.9$, $\beta_2 = 0.999$, $\epsilon = 10^{-8}$.
\end{quote}

\subsection{4. Políticas de Ajuste (\emph{Learning Rate Schedulers}) y Warmup}

\subsubsection{4.1 Planificadores de Tasa de Aprendizaje (\emph{Schedulers})}

Reducen la tasa de aprendizaje $\eta$ a medida que avanza el entrenamiento para evitar sobrepasar el mínimo local:
\begin{itemize}
	\item \textbf{Polynomial Decay:} $\eta_t = \eta_0 \left(1 - \frac{t}{T}\right)^p$.
	\item \textbf{Step / Factor Decay:} $\eta_{t+1} = \eta_t \cdot \alpha$ con $\alpha \in (0, 1)$ cada $k$ épocas.
	\item \textbf{Cosine Annealing:} Modula la tasa siguiendo una curva cosenoidal decreciente.
	\item \textbf{ReduceLROnPlateau:} Reduce $\eta$ por un factor cuando una métrica de validación deja de mejorar.
\end{itemize}

\subsubsection{4.2 Calentamiento (\emph{Warmup})}

Inicia con una tasa de aprendizaje muy pequeña durante las primeras épocas e incrementarla linealmente hasta el valor $\eta_0$. Previene la desestabilización inicial de los pesos cuando los gradientes son excesivamente volátiles.

\subsection{5. Práctica y Aplicación en Código (PyTorch)}

\begin{lstlisting}[language=Python]
	import torch
	import torch.nn as nn
	import torch.optim as optim
	from torch.optim.lr_scheduler import LambdaLR
	
	class DummyNet(nn.Module):
	def __init__(self, input_dim=10, num_classes=2):
	super(DummyNet, self).__init__()
	self.fc = nn.Sequential(
	nn.Linear(input_dim, 32),
	nn.ReLU(),
	nn.Linear(32, num_classes)
	)
	
	def forward(self, x):
	return self.fc(x)
	
	def setup_optimizers(model):
	optimizer_sgd = optim.SGD(
	model.parameters(), 
	lr=0.01, 
	momentum=0.9, 
	weight_decay=1e-4
	)
	
	optimizer_adam = optim.Adam(
	model.parameters(), 
	lr=0.001, 
	betas=(0.9, 0.999), 
	eps=1e-8
	)
	
	return optimizer_sgd, optimizer_adam
	
	def get_warmup_scheduler(optimizer, warmup_epochs=5, max_epochs=50):
	def lr_lambda(epoch):
	if epoch < warmup_epochs:
	return float(epoch + 1) / float(warmup_epochs)
	return 0.5 * (1.0 + torch.cos(torch.tensor((epoch - warmup_epochs) / (max_epochs - warmup_epochs) * 3.14159)))
	
	return LambdaLR(optimizer, lr_lambda=lr_lambda)
	
	def train_demo():
	model = DummyNet()
	optimizer = optim.Adam(model.parameters(), lr=1e-3)
	scheduler = get_warmup_scheduler(optimizer, warmup_epochs=5, max_epochs=20)
	criterion = nn.CrossEntropyLoss()
	
	inputs = torch.randn(16, 10)
	labels = torch.randint(0, 2, (16,))
	
	print("Demostraciod de Entrenamiento con Scheduler")
	for epoch in range(1, 21):
	optimizer.zero_grad()
	outputs = model(inputs)
	loss = criterion(outputs, labels)
	loss.backward()
	optimizer.step()
	
	current_lr = optimizer.param_groups[0]['lr']
	print(f"Epoca [{epoch:2d}/20] | Loss: {loss.item():.4f} | LR: {current_lr:.6f}")
	
	scheduler.step()
	
	if __name__ == "__main__":
	train_demo()
\end{lstlisting}

\section{Conexión con las prácticas}

Las prácticas comparan optimizadores y schedulers en experimentos con datos y pérdidas observables. Las curvas incluidas en las salidas permiten discutir velocidad, estabilidad y sensibilidad al learning rate.


\section{Autoevaluación razonada}

\subsection*{Pregunta} ¿Qué problema intenta aliviar momentum?
\textbf{Solución.} La oscilación y el avance lento en direcciones con curvaturas o escalas muy diferentes.

\subsection*{Pregunta} ¿Por qué un scheduler modifica el learning rate?
\textbf{Solución.} Para adaptar el tamaño de las actualizaciones durante el entrenamiento y favorecer convergencia o refinamiento final.

